% ============================================================================
% MUSTERLÖSUNG DS 1 - Gruppe A + B
% ============================================================================

\documentclass[11pt,a4paper]{article}

\usepackage[utf8]{inputenc}
\usepackage[T1]{fontenc}
\usepackage[ngerman]{babel}
\usepackage[left=2cm,right=2cm,top=1.8cm,bottom=1.5cm]{geometry}
\usepackage{helvet}
\usepackage{xcolor}
\usepackage{tabularx}
\usepackage{array}
\usepackage{enumitem}
\usepackage{tcolorbox}
\usepackage{fancyhdr}
\usepackage{lastpage}

\renewcommand{\familydefault}{\sfdefault}

\definecolor{spanisch}{HTML}{c41e3a}
\definecolor{grau}{HTML}{666666}
\definecolor{hellgrau}{HTML}{f5f5f5}
\definecolor{gruppeB}{HTML}{1565c0}
\definecolor{loesung}{HTML}{c62828}

\pagestyle{fancy}
\fancyhf{}
\fancyhead[L]{\small\textcolor{grau}{Spanisch Spätbeginner -- MUSTERLÖSUNG}}
\fancyhead[R]{\small\textcolor{grau}{ML-01}}
\fancyfoot[C]{\small\thepage/\pageref{LastPage}}
\renewcommand{\headrulewidth}{0.5pt}

\newcommand{\lsg}[1]{\textcolor{loesung}{\textbf{#1}}}
\newcommand{\es}[1]{\textit{#1}}

\newtcolorbox{gruppenbox}[1][]{
    colback=white,
    colframe=spanisch,
    fonttitle=\bfseries\color{white},
    colbacktitle=spanisch,
    title=#1,
    boxrule=1pt,
    arc=0pt,
    left=5pt, right=5pt, top=5pt, bottom=5pt
}

\newtcolorbox{gruppeBbox}[1][]{
    colback=white,
    colframe=gruppeB,
    fonttitle=\bfseries\color{white},
    colbacktitle=gruppeB,
    title=#1,
    boxrule=1pt,
    arc=0pt,
    left=5pt, right=5pt, top=5pt, bottom=5pt
}

\begin{document}

% ============================================================================
% KOPFBEREICH
% ============================================================================
\begin{center}
{\Large\textbf{\textcolor{loesung}{MUSTERLÖSUNG}}}\\[0.3em]
{\large DS 1 -- Bucharbeit (S. 30--31 / S. 102--103)}
\end{center}

\vspace{0.5em}
\hrule
\vspace{1em}

% ============================================================================
% GRUPPE A
% ============================================================================
\begin{gruppenbox}[Gruppe A: Mi barrio (S. 30--31)]

\textbf{Kasten 1: ser -- estar -- hay}

\begin{tabularx}{\textwidth}{@{}|p{2.5cm}|X|X|@{}}
\hline
\textbf{Verb} & \textbf{Verwendung} & \textbf{Beispiel} \\
\hline
\textbf{ser} & \lsg{Eigenschaften} (dauerhaft) & Mi barrio \lsg{es} grande. \\
\hline
\textbf{estar} & \lsg{Ort, Zustand} (Ort, Zustand) & El cine \lsg{está} en la calle Mayor. \\
\hline
\textbf{hay} & \lsg{Existenz} (es gibt) & \lsg{Hay} un parque y muchas tiendas. \\
\hline
\end{tabularx}

\vspace{0.8em}

\textbf{Aufgabe 1: ser / estar / hay ergänzen (6 P.)}
\begin{enumerate}[label=\alph*), leftmargin=2em, nosep]
    \item Mi barrio \lsg{es} muy tranquilo.
    \item El supermercado \lsg{está} cerca de mi casa.
    \item En mi barrio \lsg{hay} dos parques grandes.
    \item La biblioteca \lsg{es} moderna y bonita.
    \item ¿Dónde \lsg{está} el hospital?
    \item No \lsg{hay} cine en mi barrio.
\end{enumerate}

\vspace{0.5em}

\textbf{Aufgabe 2: Sätze mit vuestro (4 P.)}
\begin{enumerate}[label=\alph*), leftmargin=2em, nosep]
    \item \lsg{Vuestro barrio es grande.}
    \item \lsg{Vuestra escuela es moderna.}
    \item \lsg{Vuestros amigos son simpáticos.}
    \item \lsg{Vuestras casas son pequeñas.}
\end{enumerate}

\vspace{0.5em}

\textbf{Aufgabe 3: Wortschatz Orte (8 P.)}

\begin{tabularx}{\textwidth}{@{}|X|X|X|X|@{}}
\hline
el parque & \lsg{der Park} & la piscina & \lsg{das Schwimmbad} \\
\hline
el cine & \lsg{das Kino} & el teatro & \lsg{das Theater} \\
\hline
la biblioteca & \lsg{die Bibliothek} & el supermercado & \lsg{der Supermarkt} \\
\hline
la iglesia & \lsg{die Kirche} & el hospital & \lsg{das Krankenhaus} \\
\hline
\end{tabularx}

\vspace{0.5em}

\textbf{Aufgabe 5: Blogeintrag (8 P.) -- Beispiellösung:}

\lsg{Mi barrio se llama Stadtmitte. Es un barrio muy tranquilo y bonito. Hay un parque grande donde me gusta pasear. También hay un supermercado y una biblioteca. El cine está cerca de mi casa. Mi barrio es pequeño pero me gusta mucho porque es muy acogedor.}

\end{gruppenbox}

\vspace{1em}

% ============================================================================
% GRUPPE B
% ============================================================================
\begin{gruppeBbox}[Gruppe B: El sistema educativo (S. 102--103)]

\textbf{Kasten 2: Sistema educativo español}

\begin{tabularx}{\textwidth}{@{}|p{4cm}|c|X|@{}}
\hline
\textbf{Stufe} & \textbf{Alter} & \textbf{Beschreibung} \\
\hline
ESO & 12--16 & \lsg{Pflichtschule Sekundarstufe I} \\
\hline
Bachillerato & 16--18 & \lsg{Gymnasium / Oberstufe (2 Jahre)} \\
\hline
FP & 16+ & \lsg{Berufsausbildung} \\
\hline
\end{tabularx}

\vspace{0.8em}

\textbf{Aufgabe 1: Spanisches Schulsystem erklären (6 P.) -- Beispiel:}

\lsg{En España, los niños van a la Educación Primaria desde los 6 hasta los 12 años. Después tienen que hacer la ESO, que es obligatoria. A los 16 años pueden elegir entre el Bachillerato o la Formación Profesional.}

\vspace{0.5em}

\textbf{Aufgabe 2: Vergleich Spanien -- Deutschland (8 P.)}

Tabelle:
\begin{itemize}[leftmargin=1.5em, nosep]
    \item Notas Alemania: \lsg{1--6 oder 15--0 Punkte}
    \item Pflichtschule: \lsg{bis 16 (oder 18 in manchen Bundesländern)}
    \item Alternative zu Bachillerato: \lsg{Ausbildung / duale Berufsausbildung}
\end{itemize}

Vergleichssätze:
\begin{enumerate}[leftmargin=2em, nosep]
    \item \lsg{Mientras que en España hay un sistema centralizado, en Alemania cada Bundesland tiene su propio sistema.}
    \item \lsg{En Alemania hay Gymnasium, ya que es una escuela para estudiantes que quieren ir a la universidad.}
    \item \lsg{Además, en Alemania existe la formación dual que combina trabajo y escuela.}
\end{enumerate}

\vspace{0.5em}

\textbf{Aufgabe 3: Mediación FP-Flyer (10 P.) -- Beispiel:}

\lsg{En Alemania existe un sistema que se llama ``formación dual''. Los aprendices trabajan en una empresa y además van a una escuela profesional. El programa dura entre dos y tres años y medio. Durante este tiempo, los aprendices reciben un salario. Al final reciben un certificado que les permite trabajar en su profesión.}

\vspace{0.5em}

\textbf{Aufgabe 4: Wortschatz Bildung (6 P.)}

\begin{tabularx}{\textwidth}{@{}|X|X|X|X|@{}}
\hline
el instituto & \lsg{die Schule (Sek.)} & la formación & \lsg{die Ausbildung} \\
\hline
el bachillerato & \lsg{das Abitur / Oberstufe} & la empresa & \lsg{das Unternehmen} \\
\hline
la carrera & \lsg{das Studium / Beruf} & el aprendiz & \lsg{der Auszubildende} \\
\hline
\end{tabularx}

\end{gruppeBbox}

\vspace{1em}
\hrule
\vspace{0.3em}

\textbf{Bewertungshinweise:}
\begin{itemize}[leftmargin=1.5em, nosep]
    \item Gruppe A: 32 Punkte gesamt
    \item Gruppe B: 30 Punkte gesamt
    \item Bei Blogeintrag/Mediación: Inhalt wichtiger als perfekte Grammatik
    \item Halbe Punkte möglich bei teilweise richtigen Antworten
\end{itemize}

\end{document}
